\documentclass[11pt,a4paper]{article}
\usepackage[a4paper,margin=2.5cm]{geometry}
\usepackage{amsmath,amssymb,amsthm}
\usepackage{graphicx}
\usepackage{booktabs}
\usepackage{multirow}
\usepackage{longtable}
\usepackage{array}
\usepackage{hyperref}
\usepackage{xcolor}
\usepackage{listings}
\usepackage{enumitem}
\usepackage{parskip}
\usepackage{fancyhdr}
\usepackage{titlesec}
\sloppy
\setlength{\emergencystretch}{3em}

\hypersetup{
  colorlinks=true,
  linkcolor=blue!50!black,
  citecolor=blue!50!black,
  urlcolor=blue!50!black,
  pdftitle={ozGUI / ozLib Supplement: Multicomponent and Polydisperse Extension},
  pdfauthor={ozGUI/ozLib project}
}

\lstset{
  basicstyle=\ttfamily\small,
  breaklines=true,
  frame=single,
  columns=fullflexible,
  backgroundcolor=\color{gray!8},
  keywordstyle=\color{blue!60!black},
  commentstyle=\color{green!40!black},
  stringstyle=\color{red!60!black},
  language=Python,
  showstringspaces=false
}

\pagestyle{fancy}
\fancyhf{}
\fancyhead[L]{ozGUI / ozLib Supplement: Multicomponent Extension}
\fancyhead[R]{\thepage}
\renewcommand{\headrulewidth}{0.4pt}

% allow line breaks inside long \texttt identifiers without seqsplit,
% which cannot be used inside section titles / math
\newcommand{\code}[1]{\texttt{\hspace{0pt}#1}}
\newtheorem{remark}{Remark}

\title{\bfseries ozGUI / ozLib Supplement:\\[2mm]
Multicomponent and Polydisperse Extension\\[1mm]
\large Polydisperse hard-core Yukawa with the Rogers--Young closure,
adhesive hard spheres, and the SASfit plugin}
\author{ozGUI/ozLib project}
\date{\today}

\begin{document}
\maketitle

\begin{abstract}
\noindent
This supplement extends \emph{ozGUI / ozLib: Ornstein--Zernike Solver
Toolkit} to genuinely multicomponent systems. The scalar solver is
generalised to matrix-valued correlation functions, giving partial structure
factors $S_{ij}(q)$ for polydisperse mixtures under \emph{any} closure already
in the library. Concretely we add: a polydisperse hard-core Yukawa potential
with moment-matched size classes; the Rogers--Young closure for mixtures with
its mixing parameter fixed by thermodynamic consistency rather than supplied
by the user; two new GUI tabs; and a C implementation shipped as a SASfit
structure-factor plugin. Every claim below is accompanied by the numerical
check that supports it, and \S\ref{sec:traps} documents nine defects found
during validation, several of which produced plausible but wrong answers
rather than obvious failures.
\end{abstract}

\tableofcontents

%======================================================================
\section{Scope and design constraint}

The starting point was that the existing solver stack is strictly
one-component: $c(r)$, $h(r)$ and $\gamma(r)$ are scalar arrays of length
$N_r$, and every closure in \code{update\_c} acts on them elementwise.

The design constraint adopted here was that \textbf{one-component behaviour
must remain bit-identical}. This is verifiable rather than aspirational: a
regression suite of eight potential/closure combinations is compared against
values recorded before the change (\S\ref{sec:regression}).

A second observation made the work far smaller than expected. Every solver
class --- Picard, the two Anderson variants, the two \code{scipy} solvers,
Biggs--Andrews and both SUNDIALS KINSOL drivers --- is \emph{already}
agnostic to the shape of the unknown. Searching the whole solver layer for
\code{numberOfRadialSamplingPoints} returns no hits: each iterates a flat
vector through \code{fixPointOperator} and never inspects what it represents.
Consequently the multicomponent problem can be carried through the existing
acceleration machinery unchanged, packed as the $p(p+1)/2$ unique pairs of the
symmetric $\Gamma_{ij}(r)$.

%======================================================================
\section{Multicomponent Ornstein--Zernike equations}

\subsection{Matrix form}

For $p$ species with number densities $\rho_i$, the OZ equations couple
$p(p+1)/2$ independent pair functions,
\begin{equation}
h_{ij}(r) = c_{ij}(r)
  + \sum_{\lambda=1}^{p}\rho_\lambda \int \mathrm{d}\mathbf{r}_1\,
    h_{i\lambda}(r_1)\,c_{\lambda j}(|\mathbf{r}-\mathbf{r}_1|).
\end{equation}
In Fourier space this becomes an algebraic matrix relation. With
$\mathsf{C}(q)$ the matrix of $\hat c_{ij}(q)$ and
$\boldsymbol{\rho}=\mathrm{diag}(\rho_i)$,
\begin{equation}
\mathsf{H}(q) = \bigl[\mathsf{I}-\mathsf{C}(q)\boldsymbol{\rho}\bigr]^{-1}
                \mathsf{C}(q),
\qquad
\mathsf{\Gamma}(q) = \mathsf{H}(q)-\mathsf{C}(q).
\label{eq:ozmatrix}
\end{equation}
The single scalar division of the one-component code,
$\hat\gamma = \hat c/(1-\rho\hat c) - \hat c$, is replaced by
\eqref{eq:ozmatrix}, solved as a small dense linear system at each of the
$N_r$ grid points. Nothing else in the fixpoint operator changes: the closure
step and the two Hankel transforms are the same operations applied to
$(p,p,N_r)$ arrays.

\subsection{Implementation}

The changes to \code{oZfixpointOperator.py} and \code{oZsolver.py} are:

\begin{itemize}[nosep]
\item the Hankel transforms act on the \emph{last} axis (\code{f.shape[-1]},
      \code{dst(...,\,axis=-1)}), which is identical for 1-D input but also
      accepts the pair-matrix arrays;
\item \code{isMulticomponent()}, \code{numberOfUniquePairs()},
      \code{packPairs()} and \code{unpackPairs()} convert between the
      symmetric $(p,p,N_r)$ matrix and the flat vector the solvers expect;
\item \code{fixPointOperatorForGammaMulticomponent()} implements
      \eqref{eq:ozmatrix};
\item \code{calculateSqMulticomponent()} and a multicomponent branch in
      \code{calculateRDF()} produce the observables;
\item \code{setStartValue()} accepts the longer vector.
\end{itemize}

\begin{remark}
Because every closure in \code{update\_c} is elementwise \code{numpy}, PY,
HNC, RY, HMSA, Verlet, BPGG and the rest all become multicomponent as soon as
the potential arrays carry pair indices. Closures requiring an auxiliary
reference solve (RHNC, MHNC, RMSA, EuRah) remain one-component, as do those
reading \code{repulsivePartOfP2Ppotential} (HMSA, VM, CJVM, BB, DH, CG,
SMSA), whose split arrays are still 1-D.
\end{remark}

\subsection{Public interface: 1-D curves preserved}

The public getters deliberately continue to return \emph{one}-dimensional
curves, so the plot tabs, \code{ozLib.solve()}'s curve dictionary and the
ASCII/CSV/Excel/clipboard exports need no modification:
\begin{align}
\code{getRDF()} &\;\rightarrow\; g_{NN}(r)=\sum_{ij}x_ix_j\,g_{ij}(r),\\
\code{getSq()}  &\;\rightarrow\; S^{M}(q) \text{ (the measured structure factor)}.
\end{align}
The full matrices are exposed alongside as
\code{partialRadialDistributionFunction} and
\code{partialStructureFactor}, both of shape $(p,p,N_r)$.

%======================================================================
\section{Structure-factor conventions}
\label{sec:conventions}

Three conventions appear in this problem and confusing them is an
order-of-magnitude error, so they are stated explicitly.

\paragraph{D'Aguanno--Klein (number-fraction) form.} Used internally by the
solver, following \cite{daguanno1991}:
\begin{equation}
S^{\mathrm{DAG}}_{ij}(q) = x_i\delta_{ij} + n\,x_i x_j\,\hat h_{ij}(q).
\end{equation}

\paragraph{Ashcroft--Langreth form.} Required by
\code{PolydisperseSASBase.I\_exact}, which builds
$G_i=F_i\sqrt{\rho_i}$ and contracts $G\mathsf{S}G$:
\begin{equation}
S^{\mathrm{AL}}_{ij}(q) = \delta_{ij} + \sqrt{\rho_i\rho_j}\,\hat h_{ij}(q).
\end{equation}

\paragraph{Conversion.} The two differ by a purely kinematic factor,
\begin{equation}
\boxed{\;S^{\mathrm{AL}}_{ij} = S^{\mathrm{DAG}}_{ij}\big/\sqrt{x_i x_j}\;}
\label{eq:alconv}
\end{equation}
which is easily checked on the diagonal: $S^{\mathrm{DAG}}_{ii}/x_i
= 1+n x_i \hat h_{ii} = S^{\mathrm{AL}}_{ii}$.

\paragraph{Measured structure factor.} What a scattering experiment sees,
weighted by the form amplitudes $b_i(q)$:
\begin{equation}
S^{M}(q) = \frac{\sum_{ij} b_i b_j S_{ij}(q)}{\sum_i x_i b_i^2},
\qquad
b_i(q)\propto \sigma_i^2\,\frac{j_1(q\sigma_i/2)}{q}.
\end{equation}

Conversion \eqref{eq:alconv} is applied in exactly one place in each module.
\S\ref{sec:traps} records two separate occasions on which omitting it caused
a silent factor-of-nine error.

%======================================================================
\section{Discretising a continuous size distribution}
\label{sec:moments}

\subsection{Moment matching}

A continuous distribution $F(\sigma)$ must be reduced to $p$ classes. The
naive choice of $p$ equally spaced diameters spanning
$\langle\sigma\rangle\pm3s$ is what the existing Robertus code uses. Following
\cite{daguanno1992} we instead impose
\begin{equation}
\sum_{i=1}^{p} x_i\,\sigma_i^{\,m} = \langle\sigma^m\rangle,
\qquad m=0,1,\dots,2p-1,
\label{eq:moments}
\end{equation}
i.e.\ $2p$ equations for the $2p$ unknowns $\{\sigma_i,x_i\}$. The Schulz
distribution is a gamma distribution, so \eqref{eq:moments} is precisely
Gauss--generalised-Laguerre quadrature: with $t=1/s^2-1$, the nodes and
weights of the weight function $x^{t}e^{-x}$ give
$\sigma_i = x_i^{\mathrm{node}}\langle\sigma\rangle/(t+1)$ directly. In Python
this is \code{scipy.special.roots\_genlaguerre}; in C,
\code{gsl\_integration\_fixed\_laguerre}.

\subsection{Verification}

For $s_\sigma=0.2$ and $p=3$ the classes are
$\sigma_i=\{0.7342,\,1.0532,\,1.4527\}$ with
$x_i=\{0.2795,\,0.6304,\,0.0901\}$, and the first $2p-1=5$ moments are
reproduced to eight decimal places (Table~\ref{tab:moments}).

\begin{table}[htbp]
\centering
\caption{Moments of the discrete three-class set against the continuous
Schulz distribution, $s_\sigma=0.2$. Exact through $m=2p-1=5$.}
\label{tab:moments}
\begin{tabular}{cccl}
\toprule
$m$ & discrete & Schulz & \\
\midrule
0 & 1.00000000 & 1.00000000 & exact \\
1 & 1.00000000 & 1.00000000 & exact \\
2 & 1.04000000 & 1.04000000 & exact \\
3 & 1.12320000 & 1.12320000 & exact \\
4 & 1.25798400 & 1.25798400 & exact \\
5 & 1.45926144 & 1.45926144 & exact \\
\bottomrule
\end{tabular}
\end{table}

\subsection{How many classes are needed?}

\cite{daguanno1992} report that $p=3$ is indistinguishable from $p=5$ up to
$s_\sigma=0.3$. This was reproduced independently
(Table~\ref{tab:pconv}): going from $p=1$ to $p=3$ changes $S^M$ by $0.54$,
whereas $p=3\rightarrow5$ changes it by $0.0015$. Polydispersity is therefore
a large effect that three moment-matched classes already capture.

\begin{table}[htbp]
\centering
\caption{Convergence in the number of classes. Polydisperse hard-core Yukawa,
$\langle\sigma\rangle=250$\,\AA, $Z=200$, $n^*=0.005$, $L_B=7.01$\,\AA,
$s_\sigma=0.2$, RY with $\alpha=0.5$.}
\label{tab:pconv}
\begin{tabular}{lcccc}
\toprule
 & $g_{NN}^{\max}$ & position & $S^{M}_{\max}$ & pairs \\
\midrule
$p=1$ (monodisperse) & 2.0375 & 6.06 & 2.0356 & 1 \\
$p=3$                & 1.7368 & 6.19 & 1.7478 & 6 \\
$p=5$                & 1.7340 & 6.19 & 1.7472 & 15 \\
\midrule
\multicolumn{5}{l}{$\max|g_{NN}(1)-g_{NN}(3)|=0.3583$,\quad
                   $\max|S^M(1)-S^M(3)|=0.5407$}\\
\multicolumn{5}{l}{$\max|g_{NN}(3)-g_{NN}(5)|=0.0115$,\quad
                   $\max|S^M(3)-S^M(5)|=0.0015$}\\
\bottomrule
\end{tabular}
\end{table}

The peak sits at $r\approx6.2\langle\sigma\rangle$ against a mean spacing
$n^{-1/3}=5.85\langle\sigma\rangle$, as expected for a strongly correlated
charged system.

\begin{remark}
The moment-matching method transfers to the adhesive-sphere model, and would
let the Robertus plugin use fewer classes for the same accuracy. The specific
claim ``$p=3$ suffices'' does \emph{not} transfer: it was established for
RY--Yukawa against Monte Carlo, and sticky spheres are far more sensitive to
the contact region.
\end{remark}

%======================================================================
\section{Polydisperse hard-core Yukawa}

\subsection{Charge-parameterised form}

Following \cite{daguanno1992}, the pair potential is
\begin{equation}
\beta\varphi_{ij}(r) = A_iA_j\,\frac{e^{-\kappa r}}{r},\qquad r>\sigma_{ij},
\qquad
A_i = Z_i\sqrt{L_B}\,\frac{e^{\kappa\sigma_i/2}}{1+\kappa\sigma_i/2},
\label{eq:dkpot}
\end{equation}
with a hard core at $\sigma_{ij}=(\sigma_i+\sigma_j)/2$. The amplitude
\emph{factorises per species}, because
$\exp(\kappa\sigma_{ij}) = \exp(\kappa\sigma_i/2)\exp(\kappa\sigma_j/2)$;
this is what allows the $A_iA_j$ form.

\paragraph{Screening.} $\kappa$ is the inverse Debye length due to the
\emph{small ions}. For a salt-free suspension with monovalent counterions,
charge neutrality gives $n_{\text{counter}}=\sum_i n_iZ_i$ and hence
\begin{equation}
\kappa^2 = 4\pi L_B\sum_i n_i Z_i
\label{eq:kappa}
\end{equation}
--- the \emph{first} power of $Z$. Using $Z^2$ produces a near-ideal-gas
$S(q)$ (see \S\ref{sec:traps}).

\subsection{Size dependence of the charge}
\label{sec:charge}

How the valence scales with size is genuinely unsettled, so it is a user
input rather than a constant. \cite{daguanno1991} assume constant surface
charge density, $Z\propto\sigma^2$, and explicitly flag it as having ``little
precise experimental information''. Charge-renormalisation saturation
(Alexander \emph{et al.}) and electrophoresis on highly charged low-salt
spheres both indicate $Z\propto\sigma$ instead. Three routes are provided:

\begin{center}
\begin{tabular}{@{}p{0.30\textwidth}p{0.62\textwidth}@{}}
\toprule
\code{chargeExponent} & $Z_i = Z_{\text{ref}}(\sigma_i/\langle\sigma\rangle)^n$;
  $n=2$ constant surface charge density (default), $n=1$ linear in size,
  $n=0$ size-independent \\
\code{componentValences} & explicit per-component array, overrides the power law \\
\code{screeningValences} & a \emph{different} charge for $\kappa$, so screening
  may use the bare charge while the amplitude uses the effective one \\
\bottomrule
\end{tabular}
\end{center}

The default $n=2$ exists only so that published D'Aguanno--Klein setups
reproduce; it is not a recommendation. Charge renormalisation is deliberately
\emph{not} implemented: its definition for a polydisperse system is unsettled
(even for one component the Alexander edge-linearisation $Z_{\text{eff}}$ and
the extrapolated point charge $Q_{\text{eff}}$ disagree), and embedding one
choice inside a fitting routine would hide a contested assumption.

\paragraph{Magnitude of the effect.} Table~\ref{tab:charge} shows the choice
is comparable in size to the polydispersity effect itself, so it is not a
detail. Note also that $\alpha$-scaling and screening are coupled through
\eqref{eq:kappa}: changing the exponent moves both the amplitude prefactor
and the screening length unless \code{screeningValences} decouples them.

\begin{table}[htbp]
\centering
\caption{Effect of the charge-scaling exponent. $p=3$, $s_\sigma=0.2$,
D'Aguanno--Klein parameters, RY $\alpha=0.5$.}
\label{tab:charge}
\begin{tabular}{lcccc}
\toprule
scaling & $Z_i$ & $\kappa$ & $g_{NN}^{\max}$ & $S^M_{\max}$\\
\midrule
$n=2$ (default) & 107.8 / 221.8 / 422.0 & 0.6054 & 1.7368 & 1.7478\\
$n=1$           & 146.8 / 210.6 / 290.5 & 0.5936 & 1.9262 & 1.8381\\
$n=0$           & 200 / 200 / 200        & 0.5936 & 2.0384 & 1.7724\\
\bottomrule
\end{tabular}
\end{table}

An internal consistency check falls out of \eqref{eq:kappa}: $n=0$ and $n=1$
must share a $\kappa$, because $\sum_i x_i\sigma_i=\langle\sigma\rangle=1$ by
construction, while $n=2$ raises it by exactly
$\sqrt{\langle\sigma^2\rangle}=\sqrt{1.04}=1.0198$. Both hold to four
decimals.

\subsection{MSA-consistent \texorpdfstring{$(z,K)$}{(z,K)} form}
\label{sec:zk}

For direct comparison with the analytic MSA/RMSA path the same convention
must be used. From the polydisperse Yukawa specification,
\begin{equation}
c_{ij}(r) = K_{ij}\,\frac{e^{-z(r-\sigma_{ij})}}{r},\quad r>\sigma_{ij},
\qquad K_{ij}=K\,\delta_i\delta_j,
\qquad \delta_i = d_i\,e^{-z\sigma_i/2},
\label{eq:msapot}
\end{equation}
with $g_{ij}(r)=0$ for $r\le\sigma_{ij}$. MSA sets $c=-\beta U$, so
\begin{equation}
\beta U_{ij}(r) = -K\,\delta_i\delta_j\,\frac{e^{-z(r-\sigma_{ij})}}{r},
\end{equation}
i.e.\ \textbf{$K>0$ is attractive and $K<0$ repulsive}, and there is
\emph{no} $\sigma_{ij}$ prefactor. Note that
\code{PolydisperseOneYukawaMSA} takes $\delta_i$ \emph{directly} --- its
default \code{np.ones(N)} means $\delta_i=1$, not $d_i=1$ --- so the
exponential in \eqref{eq:msapot} must not be applied a second time
(\S\ref{sec:traps}, item 9).

%======================================================================
\section{Rogers--Young for mixtures}

\subsection{Closure}

The RY closure is applied pairwise with a \emph{single shared} mixing
function --- there is no $ij$ subscript on $f$ in \cite{daguanno1991}:
\begin{equation}
h_{ij}(r) = -1 + e^{-\beta\varphi_{ij}(r)}
  \left[1+\frac{e^{f(r)\gamma_{ij}(r)}-1}{f(r)}\right],
\qquad
f(r) = 1-e^{-\alpha r},
\label{eq:ry}
\end{equation}
interpolating between PY ($\alpha\to0$) and HNC ($\alpha\to\infty$).

\subsection{Thermodynamic consistency fixes \texorpdfstring{$\alpha$}{alpha}}
\label{sec:alpha}

$\alpha$ is not a free parameter: it is determined by requiring the
compressibility and virial routes to the pressure to agree. For a mixture this
is a \emph{scalar} condition --- which sidesteps the apparent difficulty that
a mixture has a full compressibility \emph{matrix}:
\begin{align}
\chi^{-1}_{\text{comp}} &= 1 - n\sum_{ij}x_ix_j\,\hat c_{ij}(0),
\label{eq:chicomp}\\
\chi^{-1}_{\text{vir}} &= \frac{\partial(\beta P)}{\partial n},
\label{eq:chivir}
\end{align}
with the virial pressure
\begin{equation}
\frac{\beta P}{n} = 1
  + \frac{2\pi}{3}n\sum_{ij}x_ix_j\sigma_{ij}^3 g_{ij}(\sigma_{ij}^+)
  - \frac{2\pi}{3}n\sum_{ij}x_ix_j\int r^3 g_{ij}(r)\,
    \bigl(\beta U_{ij}\bigr)'(r)\,\mathrm{d}r .
\label{eq:virial}
\end{equation}

\paragraph{The contact term.} \cite{daguanno1991}'s equation for $\beta P$
omits the second term of \eqref{eq:virial}, legitimately, because their
macroions are so strongly charged that $g_{ij}(\sigma_{ij}^+)\approx0$. That
is \emph{not} safe in general: this model has a genuine hard core and $K$ may
be small, so contact is reachable and the term is first order in the
pressure. It is included here, with $g_{ij}$ at contact obtained by linear
extrapolation from the first two grid points outside $\sigma_{ij}$. This
extrapolation is the least certain step in the whole search and is flagged as
such in the code.

\paragraph{Numerics.} Two traps, both encountered:
\begin{enumerate}[nosep]
\item the \textbf{pair potential must be held fixed} while $n$ is perturbed.
      In the charge-parameterised form \eqref{eq:dkpot} $\kappa$ depends on
      the counterion density, so rebuilding the system at $n\pm\Delta n$
      differentiates a density-dependent potential; the symptom is a residual
      stuck near $+65$ that never crosses zero. In the $(z,K)$ form
      \eqref{eq:msapot} the potential is density-independent and the problem
      does not arise.
\item \textbf{cold-start each $\alpha$}. Warm-starting across qualitatively
      different $\alpha$ destabilises the search; warm-starting across the
      $n\pm\Delta n$ perturbations is fine, and is what \cite{daguanno1991}
      do, holding $\alpha$ fixed there because it varies slowly with $n$.
\end{enumerate}

\paragraph{Result.} At $\varphi=0.2$, $K=0$, $z=2$ with three classes the
search returns $\alpha=0.2166$ with residual $1.2\times10^{-5}$ --- a clean
root, and a plausible hard-sphere value. The residual is monotone in
$\alpha$ (Table~\ref{tab:alpha}).

\begin{table}[htbp]
\centering
\caption{Thermodynamic inconsistency
$\chi^{-1}_{\text{comp}}-\chi^{-1}_{\text{vir}}$ versus $\alpha$;
$\varphi=0.2$, $z=2$, $p=3$.}
\label{tab:alpha}
\begin{tabular}{lcccc}
\toprule
$\alpha$ & 0.20 & 0.50 & 1.00 & 2.00\\
\midrule
$K=0$   & $+0.0127$ & $-0.1847$ & $-0.4002$ & $-0.6054$\\
$K=0.5$ & $-0.2219$ & $-0.3012$ & $-0.3881$ & $-0.4697$\\
\bottomrule
\end{tabular}
\end{table}

\begin{remark}[A consistent $\alpha$ need not exist]
For $K=0.5$ the residual is negative across the whole bracket: no $\alpha$
makes the two routes agree. The implementation returns the closest endpoint
and reports ``no consistent alpha found'' rather than presenting it as a
solution. This is a property of the closure, not a numerical failure.
\end{remark}

\subsection{Validation against the literature}

The one-component limit reproduces Table~1 of \cite{daguanno1991}
(their test~3: $\beta\varphi(x)=(1000/x)e^{-0.95x}$, $n\sigma^3=0.01432$).
The self-consistent $\alpha=0.3942$ gives Table~\ref{tab:dk}.

\begin{table}[htbp]
\centering
\caption{One-component RY against \cite{daguanno1991} Table~1, test~3.}
\label{tab:dk}
\begin{tabular}{lrrr}
\toprule
quantity & this work & D'Aguanno--Klein & deviation\\
\midrule
$1/S(0)$   & 93.787 & 93.80 & $0.01\%$\\
$P^{ex}$   & 36.458 & 36.45 & $0.02\%$\\
$U^{ex}$   & 20.145 & 20.09 & $0.27\%$\\
$g_{\max}$ &  2.069 &  2.08 & $0.54\%$\\
\bottomrule
\end{tabular}
\end{table}

The two quantities that \emph{enter} the consistency condition agree to
$0.01$--$0.02\%$, which is the sharpest available test. Two further checks:
\begin{itemize}[nosep]
\item \textbf{HNC limit.} At large $\alpha$ the solver gives
      $g_{\max}=1.937$, $1/S(0)=77.69$, $U^{ex}=20.56$, $P^{ex}=36.94$
      against the paper's HNC column $1.94$, $77.5$, $20.57$, $36.95$ ---
      three-digit agreement, an independent check on the closure and the
      transforms.
\item \textbf{Grid insensitivity.} Results are flat to four decimal places
      from 50 to 400 points per $\sigma$; only $1/S(0)$ shifts slightly with
      $r_{\max}$ through the $q\to0$ extrapolation. Neither paper specifies
      its grid, and the residual differences in $g_{\max}$/$U^{ex}$ are
      therefore attributable to their (unreported) $\alpha$ or $\Delta n$
      rather than to discretisation.
\end{itemize}

\subsection{Cross-validation against the analytic MSA path}

Because RY and MSA both set $c=-\beta U$ outside the core, the
\emph{derivative} $\mathrm{d}S/\mathrm{d}K$ at $K\to0$ must agree even though
the closures differ. This is the decisive test of the potential convention of
\S\ref{sec:zk}, and it is what exposed the $\delta_i$ double-counting bug.
Table~\ref{tab:msary} shows agreement after the fix.

\begin{table}[htbp]
\centering
\caption{$\mathrm{d}S_{NN}/\mathrm{d}K$ at $Q=0.5$, $z=2$, three classes.
Agreement is near-exact where both closures reduce to the same physics and
degrades with density where PY and RY genuinely differ.}
\label{tab:msary}
\begin{tabular}{lccc}
\toprule
$\varphi$ & MSA & RY & ratio\\
\midrule
0.002 & 0.03031 & 0.03085 & 1.018\\
0.020 & 0.24025 & 0.24900 & 1.036\\
0.100 & 0.41790 & 0.47266 & 1.131\\
\bottomrule
\end{tabular}
\end{table}

Before the fix the ratio was $0.14$ at all three densities, consistent with
$1/e^{-z\sigma/2}\,^2=1/0.135=7.4$.

%======================================================================
\section{Solver strategy}
\label{sec:solvers}

Three tiers are tried in order, fastest first, each strictly more robust than
the previous one.

\paragraph{Tier 1 --- Anderson (\code{KIN\_FP}) with Picard
pre-conditioning.} Fastest where it works. Anderson diverges from a cold
start on this problem: a hand-written implementation was \emph{slower} than
undamped Picard when damped enough to be stable (851 versus 677 iterations)
and diverged otherwise, and SUNDIALS' \code{KIN\_FP} diverged too, damped or
not. The remedy is a short run of plain undamped Picard --- which \emph{is}
contractive here --- before handing over. Measured: 0 pre-steps diverges, 20
already works, 40 is the time optimum; 50 is the default for margin.

\paragraph{Tier 2 --- Newton--Krylov (GMRES with \code{KIN\_LINESEARCH}).}
A different algorithm, not a variant: damped Newton steps on the residual
$F(u)=G(u)-u$ with the Jacobian applied matrix-free through GMRES. It
converged in \emph{every} regime tried, in a strikingly constant 31--40
Newton iterations regardless of difficulty --- the signature of global
convergence via linesearch. The cost is wall-clock: each Newton step requires
many GMRES inner iterations, each a full function evaluation, so it runs
3--6$\times$ slower than Anderson where Anderson works.

Tuning follows SASfit's own \code{KIN\_sasfit\_configure()}:
\code{MaxNewtonStep}$=100n$ (not $n$), \code{FuncNormTol}$=10^{-10}$ and
\code{ScaledStepTol}$=10^{-13}$ (deliberately different from each other),
\code{MaxRestarts}$=10$, \code{ETACHOICE1}.

\paragraph{Tier 3 --- Picard}, undamped then progressively damped.

\paragraph{The Python side.} \code{sundials4py} is available as a prebuilt
wheel on PyPI (version 7.8.0), so no binding generation or C++ build is
needed. Measured on the same problem, KIN\_FP is the best of the Python
solvers on every count: it agrees with the others where they converge, is
2--4$\times$ faster, and on the hard state point it fails to converge in
0.9\,s rather than converging onto a spurious root.

\begin{center}
\begin{tabular}{lccc}
\toprule
state point & KIN\_FP & \code{scipy} Anderson & Picard\\
\midrule
$\varphi=0.02$, $K=1$ & \textbf{0.1\,s} & 0.2\,s & 0.1\,s\\
$\varphi=0.1$, $K=0.2$ & \textbf{0.2\,s} & 0.8\,s & ---\\
$\varphi=0.1$, $K=1$ & rejected, 0.9\,s & rejected, 5.6\,s & diverged, 56\,s\\
\bottomrule
\end{tabular}
\end{center}

A caution recorded earlier in this project --- that
\code{sundials4py} might silently ignore \code{KINSetMAA} when it is called
after \code{KINInit}, as pyozGUI does --- was \emph{tested and found to be
unfounded}. On a deliberately slow linear fixed point, MAA${}=0$ took 757
iterations and MAA${}=5$ took 93, \emph{identical in both call orders}, so
Anderson acceleration is genuinely active. The C requirement is nonetheless
real: calling \code{KINSetMAA} after \code{KINInit} segfaults on SUNDIALS
7.1.1, so the C core keeps the documented order.

\begin{table}[htbp]
\centering
\caption{Solver tier actually used, and wall-clock time, across seven
parameter regimes (C implementation, SUNDIALS~7, $N_r=4095$).}
\label{tab:tiers}
\begin{tabular}{lcc}
\toprule
regime & time & tier\\
\midrule
$p=3$, $s=0.2$              & 0.25\,s & Anderson\\
$p=1$ monodisperse          & 0.03\,s & Anderson\\
$p=5$, $s=0.3$              & 6.85\,s & Newton--Krylov\\
$p=3$, $s=0.4$              & 0.32\,s & Anderson\\
$p=3$, $s=0.2$, $5\times$ density   & 0.29\,s & Anderson\\
$p=3$, $s=0.2$, $0.2\times$ density & 0.25\,s & Anderson\\
$p=7$, $s=0.25$             & 4.11\,s & Anderson\\
\bottomrule
\end{tabular}
\end{table}

\paragraph{Never trust the convergence flag.} \code{KINSol} can return
\code{KIN\_SUCCESS} after diverging: once the RY closure's exponent clipping
pins values near $10^{217}$ the internal stopping test is defeated. This was
observed at $p=5$, $s_\sigma=0.3$, where $S^M$ was returned as
$\sim10^{222}$ \emph{with a success flag}. Both KINSOL drivers therefore
recompute the residual after \code{KINSol} returns and reject the result
unless it really is small. With that check in place the regime correctly falls
through to Newton--Krylov and returns the right answer.

%======================================================================
\section{The SASfit plugin}

The C implementation is shipped as
\code{src/plugins/ry\_polydisperse\_yukawa/}, exporting
\code{sasfit\_sq\_RYPolydisperseYukawa}. It follows the established plugin
conventions: \code{sasfit\_cmake\_plugin()}, \code{include/private.h}, Doxygen
\code{\textbackslash defgroup}/\code{\textbackslash ingroup} with an HTML
parameter table, \code{SASFIT\_PLUGIN\_EXP\_*} registration, and the
\code{func}/\code{\_f}/\code{\_v} triad with \code{\_f} and \code{\_v} stubbed
to $0$ as in SASfit's own \code{sq\_*} plugins.

\paragraph{Transform.} The radial transform is a DST-I built from a radix-2
real FFT, which is why $N_r+1$ must be a power of two; $N_r=2^{12}-1$ at 100
points per $\sigma$ matches the Python grid exactly. The DST-I agrees with
\code{scipy.fft.dst(type=1)} exactly.

\paragraph{Caching.} A solve costs a few tenths of a second, so re-solving per
$q$-point would make any scan or fit unusable. The solved system is cached on
the full parameter set and successive solves warm-start from the previous
$\gamma$; 2000 further $q$-points after the first cost effectively nothing.

\paragraph{Portability.} Builds and runs against both SUNDIALS~7.1.1 and
6.4.1 with identical results (v7 about 30\% faster). The one relevant API
break is \code{SUNContext\_Create}, whose first argument became a typed
\code{SUNComm} with a \code{SUN\_COMM\_NULL} sentinel in v7; this is handled
by a macro switched on \code{SUNDIALS\_VERSION\_MAJOR}. Version~7 also
requires \code{-lsundials\_core}, which does not exist in v6.

\paragraph{Performance history.} Four measured steps, each kept as a test:

\begin{center}
\begin{tabular}{lr}
\toprule
change & time\\
\midrule
initial (GSL \code{LU\_decomp} per $q$-point, damped Picard) & 15.1\,s\\
inline small-matrix solve + upper-triangle symmetry          & 8.3\,s\\
undamped Picard                                              & 2.3\,s\\
Anderson (\code{KIN\_FP}) + Picard pre-conditioning          & \textbf{0.26\,s}\\
\bottomrule
\end{tabular}
\end{center}

The first version was \emph{slower than \code{numpy}}, because it performed
3.4 million GSL $3\times3$ \code{LU\_decomp} calls, where GSL's generality
dominates the actual arithmetic.

\paragraph{Agreement with Python.} The full $S^M(q)$ curve matches the Python
implementation to $7.1\times10^{-11}$ relative on the Picard path and
$1.3\times10^{-8}$ on the Anderson path (the latter being the solver
tolerance, not an error).

\paragraph{$\alpha$ in the plugin.} The consistency search costs three extra
OZ solves per trial value and is far too slow inside a fit. The plugin
therefore takes $\alpha$ as a parameter; it should be determined once in the
GUI (\S\ref{sec:alpha}) and then held fixed, which is justified because
$\alpha$ varies slowly with density --- the same reasoning
\cite{daguanno1991} use when holding it constant across their own
finite-difference density derivative.

%======================================================================
\section{New GUI tabs}

\code{oZgui.py} gained a top-level notebook and an \code{\_addExtraTabs()}
hook driven by a registry,
\begin{lstlisting}[basicstyle=\ttfamily\footnotesize]
EXTRA_TABS = [
    ("polydisperse_yukawa_tab",     "PolydisperseYukawaTab",     "Polydisperse Yukawa"),
    ("robertus_shs_tab",            "RobertusSHSTab",            "Robertus SHS"),
    ("ry_polydisperse_yukawa_tab",  "RYPolydisperseYukawaTab",   "RY Polydisperse Yukawa"),
]
\end{lstlisting}
Each tab is a self-contained \code{ttk.Frame} isolated by \code{try/except},
so adding one is a single line and an import failure costs only that tab.

Both new tabs follow the same pattern: a parameter panel; plot tabs for
$I(Q)$, the approximation error, $S_{ij}(Q)$, $S(Q)$ and the size
distribution; a Summary tab; a threaded compute handing results back through a
queue (Tk is not thread-safe); and ASCII/CSV export.

\paragraph{Approximation schemes.} Because both new SAS classes derive from
the same \code{PolydisperseSASBase}, all six of SASfit's schemes for combining
a structure factor with a size distribution --- monodisperse, decoupling
(Kotlarchyk--Chen), local monodisperse (Pedersen), partial structure factors,
scaling (Gazzillo) and van der Waals one-fluid --- come for free, and the
error of each can be read directly off the plot. The approximation code is
byte-identical between the tabs, so comparing them is a genuine test of how
those approximations fare for a different interaction rather than a comparison
of two implementations.

\paragraph{Failure reporting.} An individual scheme that fails does not lose
the run: it is marked ``unavailable'' in the summary. This matters because the
\emph{monodisperse} reference can have no PY solution at a state point where
the polydisperse system does. Adhesive-sphere failures are additionally
translated into actionable advice --- \emph{increase} $\tau$, since smaller
$\tau$ is stickier and below $\tau_c(\varphi)$ the adhesive PY closure has no
physical solution at all.

\paragraph{$\alpha$ is not exposed.} Following \S\ref{sec:alpha}, the RY tab
has no $\alpha$ entry field. It is solved for on every Compute and reported in
the Summary, together with the residual and an explicit warning when no
consistent value exists.


%======================================================================
\section{Generic polydisperse models: any potential, any closure}
\label{sec:generic}

The multicomponent machinery of \S\ref{sec:conventions} was initially reached
only through one bespoke potential. This section describes its generalisation
to \emph{any} of the library's eighteen one-component potentials, combined
with any closure that works multicomponent and any form factor.

\subsection{Reusing eighteen setters unmodified}

The potentials are not rewritten. \code{getrArray()} gains an override; the
builder points it at $r/\sigma_{ij}$ and calls the ordinary one-component
setter. Because every setter measures its hard core against
\code{hardSphereDiameter}, which is unity in those reduced units, that core
lands exactly at $\sigma_{ij}$ while the tail is evaluated at the reduced
separation. The mixing rule is therefore
\begin{equation}
\sigma_{ij}=\tfrac12(\sigma_i+\sigma_j),
\qquad
u_{ij}(r) = u\!\left(r/\sigma_{ij}\right),
\label{eq:reducedtail}
\end{equation}
i.e.\ every pair sees the same interaction \emph{shape} in units of its own
contact distance. This is the assumption Robertus \emph{et al.}\ make for
size-independent stickiness, and it is what makes the construction
potential-agnostic: $\varepsilon$, $\delta$, $n$, $\tau$ keep their
reduced-unit meaning for every pair, so no per-potential mixing rule need be
invented.

\begin{remark}
Equation \eqref{eq:reducedtail} is a modelling choice, not a derivation. For
Lennard-Jones the conventional alternative is Lorentz--Berthelot,
$\varepsilon_{ij}=\sqrt{\varepsilon_i\varepsilon_j}$ with per-species
$\varepsilon$; that is a different model and would require per-species tail
parameters, which the builder deliberately does not invent.
\end{remark}

\paragraph{Charge-coupled potentials are refused.} DLVO, DLVO-Hydra,
IonicMicrogel and the dedicated polydisperse Yukawa have amplitudes that scale
with particle size and a $\kappa$ set by the whole distribution through the
counterion density \eqref{eq:kappa}. The reduced-tail rule is simply wrong for
them, so they are rejected with a message rather than silently mis-modelled.

\paragraph{Verification.} With $s_\sigma=0$ and one class the generic path
reproduces the ordinary one-component setters \textbf{bit-identically}
(maximum difference $0.00\times10^{0}$ for hard spheres, square well and
Lennard-Jones; the first and third reproduce the recorded values $2.3561180274$
and $2.1635946756$). This is the sharpest available check that the
generalisation leaves validated physics untouched.

\subsection{Quadrature rule per distribution}
\label{sec:nodes}

All four supported distributions (Schulz, Gaussian, log-normal, Weibull) have
\emph{analytic} moments, so computing $\langle\sigma^n\rangle$ is never the
difficulty. The difficulty is the map from moments to nodes: the Golub--Welsch
Hankel eigenproblem is classically ill-conditioned even with exact input.

\begin{table}[htbp]
\centering
\caption{Condition number of the log-normal Hankel matrix built from exact
moments. Double precision carries about $10^{16}$; past that Cholesky does
\emph{not} fail, it silently returns nonsense.}
\begin{tabular}{lccc}
\toprule
 & $s=0.3$ & $s=0.4$ & $s=0.5$\\
\midrule
$N=6$  & $1.6\times10^{7}$  & $1.7\times10^{7}$  & $1.2\times10^{8}$\\
$N=8$  & $2.1\times10^{10}$ & $1.9\times10^{11}$ & $4.7\times10^{13}$\\
$N=10$ & $7.0\times10^{13}$ & $1.7\times10^{16}$ & $5.4\times10^{19}$\\
$N=12$ & $6.4\times10^{17}$ & $1.2\times10^{21}$ & $2.3\times10^{27}$\\
\bottomrule
\end{tabular}
\end{table}

A closed-form classical rule is therefore used wherever one exists, bypassing
the Hankel step entirely: generalised Gauss--Laguerre for the Schulz/gamma
weight and Gauss--Hermite for the Gaussian, both stable in double precision;
Golub--Welsch with 200-digit arithmetic only for log-normal and Weibull, which
have no classical rule. All four are moment-exact to $\sim10^{-16}$.

\paragraph{A physical guard on the Gaussian.} Gauss--Hermite places its
outermost node at $u\approx\pm2.86$, so $\sigma$ reaches zero near
$s\approx0.35$. A form factor tolerates a vanishing bin; a structure factor
does not, because a hard core must be placed at every $\sigma_{ij}$. Nodes
below $10^{-3}\langle\sigma\rangle$ are dropped and the weights renormalised
--- sacrificing moment-exactness rather than physicality. The test is
deliberately \emph{relative}: at $s=0.35$ the node sits at $10^{-4}$, positive
but still a zero-diameter particle, so a bare $\sigma>0$ check is insufficient.

\subsection{Structure factor and form factor need different resolutions}
\label{sec:nff}

The two averages behave differently in $\sigma$. $S_{ij}(Q)$ varies smoothly
with size, so three to five moment-matched classes suffice --- and the solve
costs $O(p^2)$ pair transforms, so $p$ should stay small. The form-factor
average $\langle|F|^2\rangle$ oscillates: at high $Q$ the phase $QR$ spans
about $Q\langle\sigma\rangle s_\sigma$ radians across the distribution. The
two class counts are therefore separated, with the rule of thumb
$p_{\mathrm{FF}}\gtrsim Q_{\max}\langle\sigma\rangle s_\sigma$.

The underlying defect was subtler than resolution, and is instructive. Since
$S^{\mathrm{AL}}_{ij}=\delta_{ij}+\sqrt{\rho_i\rho_j}\,h_{ij}$, interpolating
the \emph{whole} matrix onto a finer size grid smears the Kronecker delta into
a band --- which converts the incoherent sum $\sum_i|F_i|^2$ into the coherent
$|\sum_i F_i|^2$, and coherent addition of form factors is precisely what
oscillates. Only $h_{ij}$ is interpolated; the delta is re-imposed exactly.

\begin{table}[htbp]
\centering
\caption{High-$Q$ ripple against brute-force integration of
$\langle|F|^2\rangle$ with 4000 points, $s_\sigma=0.3$.}
\begin{tabular}{lc}
\toprule
 & ripple\\
\midrule
brute-force reference        & 0.038\\
5 classes (before)           & 1.91\\
$p_{\mathrm{FF}}=40$         & \textbf{0.041}\\
$p_{\mathrm{FF}}=160$        & 0.041\\
\bottomrule
\end{tabular}
\end{table}

\begin{remark}[Two plausible wrong diagnoses]
Both are recorded because each looked convincing. First, ``too few classes'':
disproved by its own test, since the ripple did not converge away even at 30
classes. Second, ``Gauss nodes are unsuited to an oscillatory integrand'': a
true statement, and replacing them with a dense grid was the right change, but
not the cause --- the ripple actually got \emph{worse} ($1.91\rightarrow3.14$).
The brute-force reference is what localised the fault, by showing
$I_{\text{dilute}}$ correct at $0.040$ while $I_{\text{exact}}$ gave $2.79$ at
$\varphi=10^{-6}$, where $S$ must be the identity. Ten lines of direct
integration should have been written before any quadrature was tuned.
\end{remark}

\subsection{Length scale and the radius convention}

The Ornstein--Zernike equations are \emph{always} solved in reduced units
(mean diameter unity): the radial grid spans only about 41 diameters, so a
physical $\sigma=100$ would place every hard core off the end of it. A mean
radius supplied by the user is applied \emph{after} the solve, to the
diameters, the number densities ($n\propto L^{-3}$, which leaves $\varphi$
invariant) and the $q$ axis ($q_{\text{red}}=QL$); the tail parameters are
already reduced and are untouched.

The scattering radius equals the hard-core radius, $R_i=\sigma_i/2$, for every
class. Verified at $R=50$: the $S(Q)$ peak moves to $Q=0.06608$ while
$Q\sigma=6.608$ (hard spheres expect $\approx2\pi$) and $QR=4.496$ at the
first $I(Q)$ minimum (exact sphere zero $4.493$) --- both unchanged from the
dimensionless case, confirming a pure change of units.

\begin{remark}[Known limitation]
Scattering size and interaction size are perfectly correlated by construction:
no independent width, no offset. That is correct for charged colloids or bare
silica, where the particle \emph{is} the hard core; it is wrong for sterically
stabilised particles whose brush is contrast-matched ($R<\sigma/2$), for
charged colloids at low salt where the effective core exceeds the physical
particle, and for solvated or partially matched shells.
\end{remark}

%======================================================================
\section{Closure suite extended and cross-validated}
\label{sec:closures}

The closure suite was checked against \code{OrnsteinZernike.jl}, an
independently written Julia library published with a peer-reviewed review.
Bridge functions were compared at prescribed $\gamma$ by evaluating
$B=\ln g+\beta u-\gamma$ from this library's $c(r)$.

\begin{table}[htbp]
\centering
\caption{Bridge functions against the independent Julia implementation.}
\begin{tabular}{lc}
\toprule
closure & $\max|B_{\text{ours}}-B_{\text{Julia}}|$\\
\midrule
Verlet                & $1.1\times10^{-16}$\\
BPGG                  & $1.1\times10^{-16}$\\
Bomont--Bretonnet     & $3.3\times10^{-16}$\\
Martynov--Sarkisov    & $2.2\times10^{-16}$\\
Charpentier--Jackse (CJVM) & identical formula\\
Zerah--Hansen (HMSA)  & identical after substitution\\
\bottomrule
\end{tabular}
\end{table}

Choudhury--Ghosh differs by parameterisation rather than error: this library
uses a density-dependent coefficient $1.0175-0.275\cdot6\varphi/\pi$ whose
zero-density limit is exactly Julia's constant $\alpha=1.01752$.

\begin{remark}
A first run reported Verlet as differing. The cause was mine: Julia's default
is $B=8/5$ and it writes the denominator as $1+B\gamma/2$, which is the same
$0.8\gamma$ used here. Checking the actual defaults before reporting a
discrepancy is the lesson; the same pattern --- an apparent disagreement that
is really an assumption of one's own --- recurs throughout this project.
\end{remark}

Three closures were added from that comparison: Khanpour
($B=\ln(1+\alpha\gamma)/\alpha-\gamma$), Modified Verlet, and Extended
Rogers--Young, whose bridge function is the ordinary RY construction with one
extra quadratic term,
\begin{equation}
\phi=\frac{e^{f\gamma}-1}{f},
\qquad
B=-\gamma+\ln\!\left(1+\phi+a\phi^2\right),
\end{equation}
so that $a=0$ reduces \emph{exactly} to RY --- verified to
$2.5\times10^{-16}$, an internal check independent of the Julia transcription.
All three inherit the multicomponent path, every solver and the GUI dropdown
automatically, because each is an elementwise expression in $\gamma$.

\begin{remark}[Naming]
This library's \code{KH} is Kovalenko--Hirata, which abbreviates identically
to Khanpour but is a different closure; the new one is therefore spelled out
in full. \code{CJVM} was confirmed to be Charpentier--Jackse, its formula
matching Julia's exactly.
\end{remark}

\paragraph{Availability by potential.} Closures needing a one-component
reference solve (Reference HNC, Modified HNC) or a hard-sphere-specific
construction (Rescaled MSA, EuRah, ZSEP) cannot be used with a polydisperse
potential, and the GUI now offers 19 of 23 rather than offering all and
failing. Martynov--Sarkisov, Vompe--Martynov and CJVM are deliberately
\emph{not} hidden: they are structurally fine and fail only because their
square-root bridge overflows $\exp(\gamma+B)$ for strongly coupled charged
systems, a limitation that applies equally in one component.

%======================================================================
\section{Consistency search for a general potential}
\label{sec:genericalpha}

The scalar condition \eqref{eq:chicomp}--\eqref{eq:virial} is available for
any potential built by the generic route. Two points differ from
\S\ref{sec:alpha}. First, $(\beta U)'$ is differentiated \emph{numerically},
because the builder reuses arbitrary setters and cannot know their analytic
derivative. Second, the reduced-tail potential \eqref{eq:reducedtail} is
density-independent, so perturbing $\varphi$ leaves it untouched --- which is
not automatic, and is exactly the trap that produced a residual stuck near
$+65$ in the charge-parameterised case.

\begin{table}[htbp]
\centering
\caption{Both routes against analytic Percus--Yevick hard spheres,
$\eta=0.3$. Residual differences are grid discretisation.}
\begin{tabular}{lcc}
\toprule
quantity & analytic & measured\\
\midrule
$\chi^{-1}$      & 10.6622 & 11.0678\\
$\beta P/\rho$   & 3.8163  & 3.8896\\
PY residual      & $+1.2482$ & $+1.3284$\\
\bottomrule
\end{tabular}
\end{table}

The Rogers--Young residual crosses zero between $\alpha=0.1$ and $0.5$ and
converges onto the HNC limit at large $\alpha$ ($-4.1495$ against HNC
$-4.1497$) --- exactly the behaviour that makes the closure work. A live
search returns $\alpha=0.2441$ with residual $-6.9\times10^{-5}$.

\begin{remark}[Judge the residual relatively]
$\chi^{-1}$ is of order 10 for a dense fluid, so an absolute residual of
$10^{-3}$ is a \emph{relative} $10^{-4}$. An absolute threshold rejected a
perfectly good root ($\alpha=0.2443$, residual $-1.2\times10^{-3}$) as ``no
consistent value found''. The criterion is now relative to $\chi^{-1}$.
\end{remark}

\begin{remark}[A false alarm worth recording]
The search was first exercised at $s_\sigma=0.2$, $p=3$, $\varphi=0.3$, where
the residual saturated near $+2.11$ with almost no $\alpha$ dependence, and
this was read as a broken virial route. With three moment-matched classes the
largest diameter reaches $\sigma\approx2.48$, so at $\varphi=0.3$ the biggest
spheres are at or past their own close packing: the \emph{state} was
unreachable, not the code wrong. Diagnosing a numerical defect from a single
pathological point, without first checking a case with a known answer, was the
error.
\end{remark}

%======================================================================
\section{Defects found during validation}
\label{sec:traps}

Each of the following produced plausible-looking wrong answers rather than an
obvious failure, which is why they are recorded.

\begin{enumerate}
\item \textbf{HMSA returned $S(q)=1$ everywhere} for hard spheres and for
every potential except Lennard-Jones, Yukawa and DLVO. The cause:
\code{repulsivePartOfP2Ppotential} and
\code{attractivePartOfP2Ppotential} were never populated for those
potentials, so $e^{-0}=1$ gave no hard-core exclusion and $c(r)=0$ became an
exact fixed point. Fixed by a general fallback in
\code{setPotentialByName()}. Independently re-confirmed afterwards: HMSA and
RY are now bit-identical for hard spheres, as the physics requires (no
attractive tail).

\item \textbf{Screening used $Z^2$ instead of $Z$.} Equation
\eqref{eq:kappa} involves the first power. The wrong version gives
$g_{\max}\approx1.006$, i.e.\ an essentially ideal gas.

\item \textbf{The virial route differentiated a density-dependent
potential} (\S\ref{sec:alpha}).

\item \textbf{\code{KINSol} reported success after diverging}
(\S\ref{sec:solvers}).

\item \textbf{\code{KINSetMAA} must precede \code{KINInit}.} The Python
binding calls them in the opposite order; doing so in C segfaults on SUNDIALS
7.1.1, because the Anderson depth is set without the corresponding arrays
being allocated. That Python tolerates it suggests the setting may simply be
ignored there --- in which case that path runs plain fixed-point iteration
rather than the Anderson acceleration it appears to request. This is worth
checking on the Python side.

\item \textbf{DST-I factor of two.} For an odd extension,
$\operatorname{Im}V[k+1]=-y[k]$; the factor 2 is already in the DFT.

\item \textbf{Warm starting silently did nothing}, because
\code{ryp\_build\_potential} cleared the stored guess --- precisely the case
in which it should be kept.

\item \textbf{A loop that exhausted its iteration budget returned success},
because the status variable was overwritten by each call to the fixpoint map.


\item \textbf{Two placeholder ``engines'' that computed nothing.} Both the
adhesive-sphere and the Rogers--Young Python cores shipped as stubs that
returned plausible numbers without solving anything.

\code{robertus\_shs\_core\_py.py}'s \code{solve()} body was a bare
\code{pass} under a docstring reading ``Dummy solve method for
compatibility'' --- the multicomponent PY $\lambda_{ij}$ system was never
formed. \code{S\_matrix()} started from \code{np.ones((p,p))} and filled
\emph{every} element, diagonal and off-diagonal alike, from a
\emph{single}-component expression $1/(1+24\varphi A)$ times an ad-hoc
``sticky correction'', clamped with \code{np.maximum(S,0.1)} to force
positivity. The size classes were therefore never coupled.

The diagnostic that exposed it is worth remembering: in the dilute limit an
Ashcroft--Langreth $\mathsf{S}$ must tend to the \emph{identity}. The
placeholder gave diagonal $\approx1$ \emph{and off-diagonal} $\approx1$, i.e.
an all-ones matrix. $I_{\text{exact}}$ then collapses from
$\sum_i\rho_iF_i^2$ to $(\sum_i\sqrt{\rho_i}F_i)^2$, inflating it by
$(\sum_i\sqrt{w_i})^2$ --- a factor depending only on the \emph{number of
classes}. Measured $I_{\text{exact}}/I_{\text{monodisperse}}$ at
$s_\sigma=0.01$, where every scheme must agree, was $1.00$, $1.97$, $2.53$,
$5.69$, $8.84$, $18.2$ for $1,2,4,8,12,24$ classes, and the dilute limit gave
$8.83$ instead of $1$.

This also explains the reported symptom that ``the exact model does not
respond to the distribution width'': polydispersity entered only through the
class weights in that spurious sum, never through the physics. The
approximation schemes call the real \code{PolydisperseSASBase} machinery and
so did respond --- hence exact and approximate disagreeing in both scale and
width dependence. Note the first diagnosis attributed this to a missing
Ashcroft--Langreth conversion; that was wrong, since the diagonal was already
correct, and only the off-diagonal test identified the true cause.

Replaced by a faithful port of the C engine (\code{py\_compute\_ab},
\code{py\_residual}, \code{build\_Q\_matrix},
\code{rshs\_structure\_matrix}). Afterwards: dilute off-diagonal
$2\times10^{-7}$, $I_{\text{exact}}/I_{\text{dilute}} = 0.99999$,
$I_{\text{exact}}/I_{\text{monodisperse}} = 1.000$ for \emph{every} class
count, $\lambda$ residual $2\times10^{-16}$, and $S(q)$ varying smoothly with
width ($0.7276$, $0.7481$, $0.8050$, $0.8760$ at
$s_\sigma=0.01,0.1,0.2,0.3$).

Separately, \code{ry\_polydisperse\_core\_py.py} had no density or volume
fraction at all in its constructor, performed no OZ iteration, and built
$S_{ij}$ from an ad-hoc closed-form blend. It was replaced by an adapter onto
the real multicomponent solver (\S\ref{sec:zk}).

\item \textbf{Unconverged and spurious solutions returned as answers.}
\code{picardIteration()} prints ``Picard did not converge after $N$ steps''
and then calls \code{derivePhysicalQuantitiesFromFixpoint()} regardless,
exposing no flag, so callers cannot tell. For an attractive Yukawa
($K=1$, $\varphi=0.1$, $z=2$) it returned a plausible $S(Q)=2.65$ after
diverging to a residual of $4.8\times10^{204}$.

Adding a residual check is necessary but \emph{not sufficient}: the RY
equations have multiple fixed points, and different solvers converge cleanly
onto different ones. At that same state point the hand-written Anderson gave
$S(0.5)=17.03$ and \code{scipy} Anderson gave $9.71$, \emph{both} with
residual ${\sim}10^{-12}$ --- and with $\min S_{NN}(q) = -9.19$ and $-17.21$.
A structure factor is a variance and cannot be negative, so both were
spurious roots. A physicality screen ($\min S_{NN} \ge 0$) is therefore
applied in addition to the residual check. Biggs--Andrews makes the same
point from the other side: it declared convergence after 21 steps and failed
the residual check.


\item \textbf{The Kronecker delta must not be interpolated.} Smearing
$\delta_{ij}$ across a finer size grid turns the incoherent $\sum_i|F_i|^2$
into the coherent $|\sum_i F_i|^2$, reinstating exactly the form-factor
oscillations that polydispersity should remove
(\S\ref{sec:nff}). It presented as a physics question --- ``at
$s_\sigma=0.3$ no oscillations should be visible'' --- and was correct as
such.

\item \textbf{A single class stored as a $(1,1,N)$ array.}
\code{isMulticomponent()} is false either way, so the scalar code path ran on
three-dimensional input and every derived curve came back shaped $(1,1,N)$.
\code{getRDF().max()} still returned the right number, which is precisely why
the bit-identical $p=1$ check passed and hid it; anything expecting a
\emph{curve} then failed with ``object too deep for desired array''.

\item \textbf{A convergence test against the wrong array.} \code{x\_0} holds
the \emph{start} value, not the converged fixpoint, so every monodisperse
reference solve was rejected. It surfaced as all six approximation schemes
reporting ``no monodisperse MSA solution'' --- a message from the base class
that looked like a physical obstruction rather than a caller's bug.

\item \textbf{Two silent rejections in \code{setPotentialByName()}.} It
counted arguments \emph{with defaults} as mandatory, making any setter that
carries optional parameters unreachable; and array-valued parameters appeared
in the GUI as numeric boxes defaulting to \code{1.0}, failing a length check.
Both returned quietly after printing, leaving the potential untouched so that
every closure produced the trivial $S(q)=1$. The array parameters are now
keyword-only, so \code{getfullargspec()[0]} --- which the GUI builds its
fields from --- lists only the numbers a user can type.

\item \textbf{GUI defaults contradicting the documented physics.} Every
potential parameter box was seeded with \code{1.0}, so a user who left
\code{chargeExponent} alone silently got $n=1$ (charge linear in size) while
the documentation promised $n=2$ (constant surface charge density) --- a
different model, with nothing on screen to say so. Fields are now seeded from
each setter's own defaults.

\item \textbf{Tk touched from a worker thread.} Reading the solver dropdown
inside the compute worker raised ``main thread is not in main loop'' and lost
the entire run. Such controls are now read on the main thread and carried into
the worker.

\item \textbf{Two convention errors, each worth about an order of
magnitude.}
\emph{(a)} The Ashcroft--Langreth conversion \eqref{eq:alconv} was omitted in
the adhesive-sphere SAS layer, inflating $I_{\text{exact}}$ by a factor that
scales with the number of classes: the ratio
$I_{\text{exact}}/I_{\text{monodisperse}}$ at $s_\sigma=0.01$, where every
scheme must agree, was $1.00$, $1.97$, $2.53$, $5.69$, $8.84$ and $18.2$ for
$1,2,4,8,12,24$ classes, and the dilute limit gave $8.83$ instead of $1$.
Because the number of classes is held fixed while sweeping the width, the
exact curve sat about nine times too high and its genuine width dependence was
visually swamped --- the reported symptom was ``the exact model looks
independent of the distribution width''.
\emph{(b)} In the RY wrapper, $\exp(-z\sigma_i/2)$ was applied to
\code{delta}, double-counting it because
\code{PolydisperseOneYukawaMSA} takes $\delta_i$ directly. RY's $K$-response
came out a factor $7.2$ too weak. Notably the \emph{first} implementation was
correct and a subsequent ``correction'' made from reading the specification
introduced the error: only the numerical cross-check of
Table~\ref{tab:msary} settled it.
\end{enumerate}

%======================================================================
\section{Regression testing}
\label{sec:regression}

The one-component regression suite covers eight combinations: PY, HNC, RY,
HMSA and MSA on hard spheres; PY and RY on Yukawa; HNC on Lennard-Jones. Each
records $g_{\max}$, $\sum g$, $c(0)$, $\sum c$, $S_{\max}$ and $S(0)$. After
the multicomponent extension all eight are \textbf{bit-identical} to the
pre-change baseline. The suite is re-run against the deployed files rather
than a working copy, and should be re-run after any change to the fixpoint
operator.

%======================================================================
\section{Known limitations}

\begin{enumerate}[nosep]
\item Closures requiring an auxiliary reference solve (RHNC, MHNC, RMSA,
      EuRah) remain one-component.
\item Closures reading the repulsive/attractive potential split (HMSA, VM,
      CJVM, BB, DH, CG, SMSA) are not yet wired for the multicomponent
      potential; those arrays remain 1-D. PY, HNC, RY, Verlet and BPGG work.
      RY is the one that matters here.
\item \code{ryp\_SM()} interpolates linearly between $q$-grid points, which
      underestimates the main peak by about $0.05\%$.
\item The plugin's cache is a single static and therefore not thread-safe.
      SASfit evaluates a model serially in $q$, which matches how it is
      driven.
\item The contact-value extrapolation in \eqref{eq:virial} is the least
      certain step of the consistency search.
\item A thermodynamically consistent $\alpha$ does not exist everywhere in
      parameter space; this is reported rather than hidden.
\item Charges are effective (renormalised), not bare (\S\ref{sec:charge}).
\end{enumerate}

%======================================================================
\begin{thebibliography}{9}

\bibitem{daguanno1991}
B.~D'Aguanno and R.~Klein,
\emph{Structural effects of polydispersity in charged colloidal dispersions},
J.~Chem.\ Soc.\ Faraday Trans.\ \textbf{87}, 379 (1991).

\bibitem{daguanno1992}
B.~D'Aguanno and R.~Klein,
\emph{Integral-equation theory of polydisperse Yukawa systems},
Phys.\ Rev.\ A \textbf{46}, 7652 (1992).

\bibitem{rogersyoung1984}
F.~J.~Rogers and D.~A.~Young,
\emph{New, thermodynamically consistent, integral equation for simple
fluids}, Phys.\ Rev.\ A \textbf{30}, 999 (1984).

\bibitem{robertus1989}
C.~Robertus, A.~P.~Philipse, J.~G.~H.~Joosten and Y.~K.~Levine,
\emph{Solution of the Percus--Yevick approximation of the multicomponent
adhesive spheres system applied to the small angle X-ray scattering from
microemulsions}, J.~Chem.\ Phys.\ \textbf{90}, 4482 (1989).

\bibitem{blumhoye1978}
L.~Blum and J.~S.~H\o ye,
\emph{Solution of the Ornstein--Zernike equation with Yukawa closure for a
mixture}, J.~Stat.\ Phys.\ \textbf{19}, 317 (1978).

\bibitem{blumarias2006}
L.~Blum and M.~Arias,
\emph{Structure of multi-component/multi-Yukawa mixtures},
arXiv:cond-mat/0602477 (2006).

\bibitem{wagner1991}
N.~J.~Wagner, R.~Krause, A.~R.~Rennie, B.~D'Aguanno and J.~Goodwin,
\emph{The microstructure of polydisperse, charged colloidal suspensions by
light and neutron scattering}, J.~Chem.\ Phys.\ \textbf{95}, 494 (1991).

\bibitem{rastogi1996}
S.~R.~Rastogi, N.~J.~Wagner and S.~R.~Lustig,
\emph{Rheology, self-diffusion, and microstructure of charged colloids under
simple shear by massively parallel dynamic simulations},
J.~Chem.\ Phys.\ \textbf{104}, 9249 (1996).

\bibitem{lee1995}
L.~L.~Lee,
\emph{An accurate integral equation theory for hard spheres: role of the
zero-separation theorems in the closure relation},
J.~Chem.\ Phys.\ \textbf{103}, 9388 (1995).

\end{thebibliography}

\end{document}
